\section{研究计划进度安排及预期目标}

\subsection{研究计划进度安排}

本研究计划周期为 12 个月，围绕基于激光雷达数据的雪深反演方法展开系统研究。研究过程综合理论分析、数据处理与模型构建，整体进度安排如下：

\textbf{第一阶段：文献调研与理论基础构建（第 1--2 个月）}

系统查阅并梳理国内外关于激光雷达测量雪深、雪面高度反演、激光回波特性分析及相关反演算法的研究文献，重点关注地基激光雷达、机载与星载激光高度计在雪深测量中的应用方法及其优缺点。在此基础上，总结激光雷达测距原理与雪深反演的理论基础，明确雪深与雪面高度、地表高度之间的基本关系：
\begin{equation}
H_s = h_{\text{snow}} - h_{\text{ground}},
\end{equation}
其中 $H_s$ 表示雪深，$h_{\text{snow}}$ 与 $h_{\text{ground}}$ 分别表示雪面高度和地表高度。通过该阶段研究，明确本研究的技术路线、关键科学问题与研究目标，完成开题报告的撰写与论证。

\textbf{第二阶段：数据获取与预处理方法研究（第 3--4 个月）}

根据研究目标，获取或整理用于雪深反演的激光雷达数据，包括实验获取的数据或已有公开数据集。针对激光雷达点云或回波数据的特点，研究数据预处理方法，重点包括噪声点剔除、异常回波筛选以及雪面与地表点的初步区分等关键步骤，为后续雪深反演模型的构建提供高质量的数据基础。

\textbf{第三阶段：雪深反演模型构建与算法实现（第 5--7 个月）}

在分析激光雷达回波特征与雪深物理关系的基础上，构建雪深反演模型。结合激光雷达测距基本原理：
\begin{equation}
R = \frac{c \, \Delta t}{2},
\end{equation}
其中 $R$ 为激光测距结果，$c$ 为光速，$\Delta t$ 为激光往返传播时间，将测距信息与系统几何模型相结合，实现雪面三维坐标及高度信息的反演。在此基础上，采用几何建模方法、统计分析方法或反演算法计算雪深信息，并在必要时引入多参数修正或辅助约束条件，以提高反演结果的精度与稳定性。通过仿真或实验数据对所构建模型进行实现与初步验证。

\textbf{第四阶段：结果分析与模型优化（第 8--9 个月）}

对雪深反演结果进行系统分析，评估模型在不同地形条件、观测角度及数据质量下的反演精度与误差特性，并与参考数据或已有方法进行对比验证。针对地形起伏、激光入射角变化及回波不确定性等主要误差来源，分析其对雪深反演结果的影响规律，并对模型进行针对性的优化与改进。

\textbf{第五阶段：论文撰写与成果总结（第 10--12 个月）}

在前期研究工作的基础上，系统整理理论分析、模型构建及实验结果，完成毕业论文的整体撰写，包括研究背景、方法原理、实验结果与讨论等内容。对论文进行修改与完善，形成结构完整、逻辑严谨、论证充分的毕业论文，并完成答辩准备工作。

\subsection{预期研究目标}

通过本研究，预期达到以下研究目标：

\textbf{理论目标：}
从激光雷达测距原理与空间几何关系出发，系统分析激光雷达回波特性与雪深之间的物理与数学联系，明确雪面高度反演在雪深测量中的理论基础。

\textbf{方法目标：}
构建一套基于激光雷达数据的雪深反演方法或模型，实现由激光雷达观测信息到雪深参数的定量反演过程，并保证该方法在典型应用场景下具有良好的可行性与稳定性。

\textbf{技术目标：}
完成雪深反演算法的实现与验证，形成较为完整的数据处理与分析流程，包括数据预处理、雪面提取、雪深计算及结果评估等关键环节。

\textbf{应用目标：}
验证所提出雪深反演方法在实际观测数据或仿真数据中的有效性，为激光雷达在雪深监测、冰雪环境研究及相关工程应用中的进一步推广提供技术参考。

\textbf{成果目标：}
完成一篇结构完整、内容规范、论证严谨的毕业论文，相关研究成果可为后续深入研究或工程应用提供基础与参考。
